\diarytitle{0111}{有限区間における複素指数関数の積分とsinc関数}

オイラーの公式 $e^{i\theta} = \cos\theta + i\sin\theta$ を用いて複素指数関数の積分を計算し、実部・虚部を比較することで実積分の値を求める。

$k \neq 0$ のとき $\lambda = -ik \neq 0$ として原始関数を求めると
\[
\int_{-a}^{a} e^{-ikx}\,dx
= \left[\frac{e^{-ikx}}{-ik}\right]_{-a}^{a}
= \frac{e^{-ika} - e^{ika}}{-ik}
\]
オイラーの公式より $e^{-ika} = \cos(ka) - i\sin(ka)$, $e^{ika} = \cos(ka) + i\sin(ka)$ であるから
\[
e^{-ika} - e^{ika} = -2i\sin(ka)
\]
よって
\[
\int_{-a}^{a} e^{-ikx}\,dx = \frac{-2i\sin(ka)}{-ik} = \frac{2\sin(ka)}{k}
\]
これを $2a$ で括り直すと
\[
\frac{2\sin(ka)}{k} = 2a\cdot\frac{\sin(ka)}{ka}
\]
したがって
\[
\boxed{\; \int_{-a}^{a} e^{-ikx}\,dx = 2a\,\mathrm{sinc}(ka) \;}
\]

（検算: 被積分関数の実部 $\cos(kx)$ は偶関数なので $\displaystyle\int_{-a}^{a}\cos(kx)\,dx = 2\int_{0}^{a}\cos(kx)\,dx = \dfrac{2\sin(ka)}{k}$ となり実部が一致し、虚部 $-\sin(kx)$ は奇関数なので積分は $0$ となって結果が実数であることとも整合する。また $k \to 0$ の極限では $\mathrm{sinc}(0)=1$ より積分値は $2a$ に収束し、これは $\displaystyle\int_{-a}^{a}1\,dx = 2a$ と一致する。）
